\documentclass[10pt, a4paper]{scrartcl}

\usepackage{vorschule}
\usepackage[
    typ=ab,
    fach=Mathematik,
    lerngruppe={6},
    nummer=6,
    module={Symbole,Lizenzen,Papiertypen},
    seitenzahlen=keine,
    farbig,
    lizenz=cc-by-nc-sa-4,
]{schule}

\usepackage[
	kuerzel=Ngb,
	reihe={Brüche},
	version={2019-10-9},
]{ngbschule}

\author{J. Neugebauer}
\title{Mit Brüchen rechnen}
\date{\Heute}

\setzeAufgabentemplate{schule-binnen}


\begin{document}

\ReiheTitel

\medskip
\begin{minipage}{0.7\textwidth}
Anne hat drei Geschwister, die fünf, sechs und acht Jahre alt sind. Als sie 
alleine ihre Oma besucht steckt die ihr einen Geldschein zu und sagt: \medskip

\enquote{\textit{Schön das du heute da warst, Anne. Das Geld ist für dich und deine 
Geschwister. Ich möchte das du jedem den Anteil abgibst, der seinem Alter entspricht.}} \medskip

\enquote{\textit{Och Oma}}, ruft Anne entsetzt, \enquote{\textit{dann bekomme ich ja am wenigsten von allen!}}
\end{minipage}\hfill
\begin{minipage}{4cm}
	\includegraphics[width=4cm]{6-AB.6-Abb_1}
\end{minipage}

\begin{center}\footnotesize
Falls ihr nicht weiter kommt:\\
Zu jeder Aufgabe gibt es vorne einen Tipp mit der entsprechenden Aufgabennummer.
\end{center}

\begin{aufgabe}
	Prüft mit den \textbf{Bruchsteinen} Annes Befürchtung. Hat sie Recht?
\end{aufgabe}
\medskip

\begin{aufgabe}
	Welchen Anteil muss Anne insgesamt abgeben? \textbf{Notiert eine Rechnung:}
	
	\feldLin{1}
\end{aufgabe}
\medskip

\begin{aufgabe}
	\enquote{Haben Brüche denselben Nenner, dann kann man sie leicht addieren und subtrahieren.}
	
	Prüft diese Aussage mit verschiedenen \textbf{Kombinationen von Bruchsteinen}. Notiert jeweils eine Rechnung und das Ergebnis als Bruch (z.B. $\frac{1}{4} + \frac{1}{4} = $...).
	
	\feldLin{4}
\end{aufgabe}
\medskip

\begin{aufgabe}
	Formuliert eine Regel, wie man Brüche mit gleichem Nenner (gleichnamige Brüche) addieren und subtrahieren kann.
	
	Man kann gleichnamige Brüche addieren und subtrahieren, indem man ...
	
	\feldLin{6}
\end{aufgabe}

\end{document}
